\documentclass[12pt,a4paper]{article}

% --- PACKAGES DE BASE ---
\usepackage[utf8]{inputenc}
\usepackage[T1]{fontenc}
\usepackage[french]{babel}
\usepackage{geometry}
\geometry{top=2.5cm, bottom=2.5cm, left=2.5cm, right=2.5cm}
\usepackage{graphicx}
\usepackage{booktabs}
\usepackage{xcolor}
\usepackage{hyperref}
\usepackage{listings}
\usepackage{fancyhdr}
\usepackage{amsmath}

% --- CONFIGURATION GRAPHIQUE & COULEURS ---
\definecolor{primary}{HTML}{1D4ED8}     % Bleu primaire
\definecolor{secondary}{HTML}{0A2647}   % Bleu nuit sidebar
\definecolor{darkslate}{HTML}{1E293B}   % Slate 800
\definecolor{lightbg}{HTML}{F8FAFC}     % Slate 50
\definecolor{bordercolor}{HTML}{E2E8F0} % Slate 200

\hypersetup{
    colorlinks=true,
    linkcolor=primary,
    filecolor=primary,      
    urlcolor=primary,
    pdftitle={IncidentFlow - Plan de Developpement},
    pdfpagemode=FullScreen,
}

% --- SETUP LISTINGS POUR CONFIG/CODE ---
\lstset{
    basicstyle=\ttfamily\small,
    backgroundcolor=\color{lightbg},
    frame=single,
    rulecolor=\color{bordercolor},
    breaklines=true,
    postbreak=\mbox{\textcolor{red}{$\hookrightarrow$}\space},
    showstringspaces=false
}

% --- EN-TÊTES ET PIEDS DE PAGE ---
\pagestyle{fancy}
\fancyhf{}
\lhead{\textcolor{secondary}{\textbf{IncidentFlow}}}
\rhead{Plan de Développement}
\cfoot{\thepage}
\renewcommand{\headrulewidth}{0.4pt}

% --- INFORMATIONS DU DOCUMENT ---
\title{
    \vspace{2cm}
    \textbf{\Huge IncidentFlow}\\[0.5cm]
    \Large Plateforme Dynamique de Gestion d'Incidents\\[1.5cm]
    \textbf{\Large Rapport de Spécification et Plan de Développement}
    \vspace{1.5cm}
}
\author{\textbf{Anas HADDOU} \\ Élève Ingénieur en Génie Informatique}
\date{\today}

\begin{document}

% Page de garde
\maketitle
\thispagestyle{empty}
\clearpage

% Table des matières
\tableofcontents
\clearpage

% --- SECTION 1 : INTRODUCTION ---
\section{Introduction}
Le projet \textbf{IncidentFlow} vise à concevoir et à implémenter une plateforme centralisée et dynamique pour la gestion du cycle de vie des incidents au sein d'une organisation. 

Contrairement aux outils traditionnels rigides, la valeur ajoutée d'IncidentFlow réside dans son \textbf{moteur de workflow dynamique}. Ce moteur permet aux administrateurs de configurer librement les états et les transitions autorisés pour chaque catégorie d'incidents directement depuis l'interface utilisateur, sans nécessiter de modification du code source.

Ce document fournit le plan de développement pas à pas ainsi que les estimations de temps et de ressources nécessaires à la réalisation de l'application, en se basant sur le document de spécification \texttt{IncidentFlow.pdf} et sur l'identité visuelle définie dans la maquette \texttt{wireframe.html}.

% --- SECTION 2 : ARCHITECTURE TECHNIQUE ---
\section{Architecture Technique et Infrastructure}
L'application adopte une architecture micro-services conteneurisée à l'aide de Docker. Elle est articulée autour de quatre conteneurs principaux décrits dans le tableau~\ref{tab:architecture} :

\begin{table}[htbp]
    \centering
    \caption{Description des conteneurs de l'infrastructure}
    \label{tab:architecture}
    \vspace{0.3cm}
    \begin{tabular}{llll}
        \toprule
        \textbf{Conteneur} & \textbf{Image de base} & \textbf{Port exposé} & \textbf{Rôle} \\
        \midrule
        \texttt{incidentflow-postgres} & PostgreSQL 15-alpine & 5432 & Persistance des données relationnelles \\
        \texttt{incidentflow-keycloak} & Keycloak (Latest) & 8180 & Gestion des identités et accès (IAM, OAuth2) \\
        \texttt{incidentflow-backend} & Spring Boot (JDK 17) & 8080 & API REST et moteur de règles métier \\
        \texttt{incidentflow-frontend} & React (Vite / Node) & 3000 & Interface utilisateur Web réactive \\
        \bottomrule
    \end{tabular}
\end{table}

\subsection{Modèle de Données Conceptuel}
Le schéma relationnel PostgreSQL s'organise autour des entités suivantes :
\begin{itemize}
    \item \textbf{Role} : Identifiant et libellé des rôles système (\textit{Administrateur}, \textit{Opérateur}, \textit{Opérateur médical}, \textit{Responsable}).
    \item \textbf{User} : Comptes utilisateurs reliés à Keycloak par leur identifiant unique.
    \item \textbf{Incident} : Informations de l'incident (titre, description, catégorie, priorité, état actuel, créateur et assignataire).
    \item \textbf{Comment \& Attachment} : Commentaires et pièces jointes associés à un incident.
    \item \textbf{Workflow, State \& Transition} : Entités de configuration dynamique des processus.
\end{itemize}

% --- SECTION 3 : STYLE ET CHARTE GRAPHIQUE ---
\section{Spécifications Graphiques (Design System)}
L'interface de l'application respectera strictement les tokens CSS définis dans le fichier \texttt{wireframe.html} pour offrir un rendu professionnel et haut de gamme.

\subsection{Palette de Couleurs}
\begin{itemize}
    \item \textbf{Arrière-plan principal} : \texttt{\#F8FAFC} (Slate 50)
    \item \textbf{Barre latérale (Sidebar)} : \texttt{\#0A2647} (Bleu nuit profond)
    \item \textbf{Couleur primaire (Actions \& Liens)} : \texttt{\#3B82F6} (Bleu Azure / Primary 500)
    \item \textbf{Textes} : \texttt{\#1E293B} (Texte principal) et \texttt{\#64748B} (Texte secondaire/muet)
\end{itemize}

\subsection{Typographie}
Le rendu textuel s'appuiera sur deux polices Google Fonts :
\begin{enumerate}
    \item \textbf{\textit{Plus Jakarta Sans}} : Utilisée pour l'ensemble des éléments de l'interface utilisateur (titres, boutons, formulaires).
    \item \textbf{\textit{JetBrains Mono}} : Réservée aux codes techniques et identifiants d'incidents (ex. : \texttt{INC-2026-001}).
\end{enumerate}

% --- SECTION 4 : PLAN DE DEVELOPPEMENT ---
\section{Plan de Développement Étape par Étape}
Le développement est découpé en cinq phases séquentielles afin de garantir une intégration et une validation progressives.

\subsection{Phase 1 : Initialisation de l'Infrastructure et de la Base de Données}
\begin{enumerate}
    \item Écriture du fichier \texttt{docker-compose.yml} configurant les volumes réseau et l'ordre de démarrage (\texttt{PostgreSQL} $\rightarrow$ \texttt{Keycloak} $\rightarrow$ \texttt{Backend} $\rightarrow$ \texttt{Frontend}).
    \item Création du Realm et du client OpenID Connect dans Keycloak, suivi de la définition des rôles applicatifs.
    \item Initialisation de la base de données PostgreSQL et écriture des scripts d'initialisation de schéma (\texttt{schema.sql}).
\end{enumerate}

\subsection{Phase 2 : Développement du Backend (Spring Boot API)}
\begin{enumerate}
    \item Sécurisation des endpoints de l'API via Spring Security en configurant la validation des tokens JWT Keycloak.
    \item Développement des services d'accès aux données (JPA Repositories) et des contrôleurs REST pour la gestion des incidents.
    \item Implémentation du moteur de validation des transitions : vérification du rôle requis, de la présence d'un commentaire obligatoire et historique automatique des changements d'état.
\end{enumerate}

\subsection{Phase 3 : Développement du Frontend (React UI)}
\begin{enumerate}
    \item Initialisation du projet React avec Vite et intégration du fichier CSS global contenant le design system (\texttt{index.css}).
    \item Intégration de la bibliothèque d'authentification Keycloak JS pour gérer la session utilisateur.
    \item Conception des composants structurels fixes (Sidebar de \texttt{260px}, barre supérieure avec indicateur de profil et centre de notifications).
    \item Création de la page d'accueil (Dashboard) avec ses cartes de métriques et listes d'incidents filtrables.
\end{enumerate}

\subsection{Phase 4 : Gestion des Fiches Incidents et Collaboration}
\begin{enumerate}
    \item Formulaire de déclaration d'incident avec upload de fichiers.
    \item Fiche détaillée de l'incident affichant les détails à gauche et le volet d'action collaborative à droite (timeline historique, commentaires et pièces jointes).
    \item Rendu dynamique des boutons de transition d'état sur la fiche selon les droits de l'utilisateur.
\end{enumerate}

\subsection{Phase 5 : Administration des Workflows et Tests E2E}
\begin{enumerate}
    \item Création du panneau d'administration permettant de modifier visuellement les états et les transitions associés à chaque catégorie d'incidents.
    \item Système d'envoi et de réception des notifications applicatives en temps réel.
    \item Écriture et passage des scénarios de test bout en bout (recette fonctionnelle).
\end{enumerate}

% --- SECTION 5 : ESTIMATIONS ---
\section{Estimations de Charge et de Ressources}
Le tableau~\ref{tab:estimations} présente l'estimation de la charge de travail pour chaque phase, comparée entre un développeur humain classique et le processus automatisé par l'agent IA Antigravity.

\begin{table}[htbp]
    \centering
    \caption{Estimation de la charge de travail par phase}
    \label{tab:estimations}
    \vspace{0.3cm}
    \begin{tabular}{lll}
        \toprule
        \textbf{Phase} & \textbf{Développeur Humain} & \textbf{Agent IA (Antigravity)} \\
        \midrule
        Phase 1 : Infrastructure \& DB & 1 à 2 jours & 20 à 30 minutes \\
        Phase 2 : Backend \& API REST & 3 à 5 jours & 45 à 60 minutes \\
        Phase 3 : Frontend \& UI de Base & 3 à 4 jours & 45 à 60 minutes \\
        Phase 4 : Fiche Incident \& Collaboration & 2 à 3 jours & 30 à 45 minutes \\
        Phase 5 : Administration \& Recette & 2 à 3 jours & 30 à 45 minutes \\
        \midrule
        \textbf{Total Estimé} & \textbf{11 à 17 jours} & \textbf{2,5 à 4 heures} \\
        \bottomrule
    \end{tabular}
\end{table}

\subsection{Consommation de Tokens (Modèle IA)}
Pour l'exécution automatisée avec le modèle \textit{Gemini 3.5 Flash (Medium)}, les volumes de tokens projetés sont :
\begin{itemize}
    \item \textbf{Tokens d'entrée (Context Window)} : Entre $400\,000$ et $750\,000$ tokens au total (itérations et relectures de fichiers incluses).
    \item \textbf{Tokens de sortie (Génération de code)} : Entre $80\,000$ et $150\,000$ tokens de code source rédigé.
\end{itemize}

\end{document}
